\section{Proof of the FLP Theorem}\label{sec:proof}

\subsection{Setup and Key Definitions}

\begin{smjdefinition}[Configuration]\label{def:flp-configuration}
A \emph{configuration} $C$ consists of the local state of every process together with the set of messages in transit. An initial configuration has every process in its initial state, determined by its input value, and has no messages in transit. A transition consists of one atomic step by a process: the process either receives one pending message or receives no message (a null receive), then performs its deterministic local computation and sends any resulting messages. A decision, once made, is irrevocable. This instantiates the generic notion of a configuration from \Cref{def:transition-system}: the global state $C$ is, concretely, the pair (local process states, messages in transit), and the transition relation is given by delivering one message at a time.
\end{smjdefinition}

\begin{smjdefinition}[Admissible Execution]
An execution is \emph{admissible} if at most one process crashes and every message sent to a non-faulty process is eventually delivered. The scheduler may choose the delivery order and may delay messages by arbitrarily large finite amounts.
\end{smjdefinition}

\begin{smjdefinition}[Valency]
A configuration $C$ is \emph{$v$-valent} (for $v\in\{0,1\}$) if every admissible execution from $C$ eventually decides $v$. It is \emph{bivalent} if there exist admissible executions from $C$ that decide $0$ and $1$. Under the assumed termination property, every configuration is therefore either 0-valent, 1-valent, or bivalent.
\end{smjdefinition}

Intuitively, a bivalent configuration is one from which the scheduler can still steer the execution toward either decision value; the algorithm has not yet committed to a unique outcome.

\subsection{The Proof}

\vspace{18pt}

\begin{proof}[Proof of \Cref{thm:flp}]
Assume for contradiction that a deterministic consensus algorithm $\mathcal{A}$ exists that tolerates one crash failure in an asynchronous system.
\\[20pt]
\textbf{Step 1: There exists a bivalent initial configuration.}
\\

Suppose for contradiction that every initial configuration is univalent. Consider the set of all possible initial configurations parameterised by process input values $v_1,\ldots,v_n\in\{0,1\}$. By validity, the all-0 configuration is 0-valent and the all-1 configuration is 1-valent. Changing inputs one process at a time, there must be two adjacent configurations $C_0$ (0-valent) and $C_1$ (1-valent) that differ only in the input of a single process $p$.

Now consider an admissible execution in which $p$ crashes before taking any step. By the assumed crash tolerance and termination property, the remaining processes eventually decide. Their local states and received messages are identical whether the initial configuration was $C_0$ or $C_1$: their own inputs are identical in both configurations, and $p$ is silent from the start. Because $\mathcal{A}$ is deterministic, the same sequence of steps by the surviving processes must therefore lead to the same decision in both executions. But $C_0$ is 0-valent and $C_1$ is 1-valent, so such a common decision is impossible. Hence our assumption was false, and \emph{some initial configuration is bivalent}.
\\[12pt]
\textbf{Step 2: From any bivalent configuration, a bivalent configuration is reachable after any applicable event.}
\\

Let $C$ be a bivalent configuration and let $e=(p,m)$ be an event applicable at $C$. Let $\mathcal{E}$ denote the set of configurations reachable from $C$ without applying $e$; $e$ remains applicable throughout $\mathcal{E}$ because messages may be delayed arbitrarily and applying other events does not remove $m$. Let $\mathcal{D}=e(\mathcal{E})$ be the set of configurations obtained by applying $e$ to configurations in $\mathcal{E}$.

We repeatedly use the basic fact that \emph{valency is preserved forward}: if a configuration is $v$-valent, then every configuration reachable from it is also $v$-valent. This follows because every admissible execution from the successor can be extended to an admissible execution from the original configuration.

Suppose, for contradiction, that every configuration in $\mathcal{D}$ is univalent. Because $C$ is bivalent, there exists a 0-valent configuration and a 1-valent configuration reachable from $C$. For each $i\in\{0,1\}$, choose an $i$-valent configuration $E_i$ reachable from $C$. If $E_i\in\mathcal{E}$, then $F_i=e(E_i)\in\mathcal{D}$ is univalent and is reachable from $E_i$, so forward preservation implies that $F_i$ is $i$-valent. Otherwise, $e$ was applied somewhere on the path from $C$ to $E_i$. Let $F_i\in\mathcal{D}$ be the configuration immediately after the first application of $e$ on that path. Then $E_i$ is reachable from $F_i$, and since $F_i$ is univalent it must again be $i$-valent. Thus $\mathcal{D}$ contains both a 0-valent and a 1-valent configuration.

Call two configurations \emph{neighbours} if one results from the other in a single event. By considering a finite path in $\mathcal{E}$ from a configuration leading to a 0-valent member of $\mathcal{D}$ to one leading to a 1-valent member, there exist neighbouring configurations $C_0,C_1\in\mathcal{E}$ such that
\[
D_i=e(C_i)
\]
is $i$-valent for $i\in\{0,1\}$. Without loss of generality, let $C_1=e'(C_0)$, where $e'=(p',m')$.

\emph{Case A: $p'\ne p$.} The events $e$ and $e'$ affect different processes and therefore commute:
\[
e(e'(C_0))=e'(e(C_0)).
\]
Hence
\[
D_1=e'(D_0).
\]
But $D_0$ is 0-valent, so every configuration reachable from $D_0$, including $D_1$, must also be 0-valent. This contradicts the fact that $D_1$ is 1-valent.

\emph{Case B: $p'=p$.} Consider a finite deciding execution from $C_0$ in which process $p$ takes no steps; such an execution exists by the assumed tolerance of one crash failure and termination. Let $\sigma$ be its event sequence and let $A=\sigma(C_0)$. Because $\sigma$ contains no event of $p$, the sequence $\sigma$ can also be applied after either $D_0=e(C_0)$ or $D_1=e(C_1)$. Since $D_0$ is 0-valent and $D_1$ is 1-valent,
\[
E_0=\sigma(D_0)
\]
is 0-valent and
\[
E_1=\sigma(D_1)
\]
is 1-valent. Moreover, by commutativity with the events of $\sigma$,
\[
E_0=e(A)
\]
and, because $C_1=e'(C_0)$,
\[
E_1=e(e'(A)).
\]
Thus from $A$ one can reach both a 0-valent and a 1-valent configuration. Hence $A$ is bivalent. But $A$ is reached by a deciding execution from $C_0$, so $A$ already has a decision value and is therefore univalent. This is a contradiction.

In both cases we obtain a contradiction. Therefore $\mathcal{D}$ contains a bivalent configuration. Equivalently, after any applicable event $e$, there is a finite execution fragment ending with $e$ whose final configuration is bivalent. This is the central bivalency-preservation lemma of Fischer, Lynch, and Paterson~\cite{flp1985}.
\\[20pt]
\textbf{Step 3: Constructing an infinite non-deciding execution.}
\\

Start from the bivalent initial configuration guaranteed by Step~1. Maintain a fixed fair ordering of the processes, and at each round select the next process $p$ in that ordering. Let $m$ be the oldest message currently pending for $p$, if one exists; otherwise let $m=\bot$ denote a null receive event in which $p$ takes a step without receiving a message. Let $e=(p,m)$.

By Step~2, there is a finite execution fragment that avoids applying $e$ until its final event and ends with a bivalent configuration after applying $e$. Use this fragment as the current round. Since $e$ is the final event of the round, the selected process makes progress at every round, while the final configuration remains bivalent.

Repeat this construction indefinitely, cycling fairly through all processes. A process is selected infinitely often, and whenever a message is pending for that process, the oldest such message is eventually selected and delivered. Hence every message sent to a non-faulty process is eventually delivered. The construction contains no crashed process at all, so the resulting infinite execution is admissible.

Because every round ends in a bivalent configuration, no process can have reached an irrevocable decision: once a decision occurs, the configuration is univalent. Thus the constructed admissible execution never reaches a decision, contradicting the assumed termination property of $\mathcal{A}$. Therefore no such deterministic consensus algorithm exists.
\end{proof}

% ── FIGURE: illustrating Steps 1 and 2 of the proof ──
\begin{figure}[H]
\centering
\begin{subfigure}[t]{\textwidth}
\centering
\begin{tikzpicture}[
    config/.style={draw, thick, minimum width=2.0cm, minimum height=0.9cm,
                   font=\small, align=center},
    zval/.style={config, fill=blue!10, draw=blue!60},
    oval/.style={config, fill=red!10, draw=red!60},
    bival/.style={config, fill=yellow!20, draw=orange},
    arrow/.style={-{Stealth[length=6pt]}, thick}
]

\node[bival] (C) at (0, 0) {$C$ (bivalent)};
\node[zval] (C0) at (-3.6, -2.4) {$e(C_0)$ (0-valent)};
\node[oval] (C1) at (1.0, -2.4) {$e(C_1)$ (1-valent)};
\node[bival] (B) at (-1.3, -4.1) {$B$ (boundary in $\mathcal{E}$)};
\node[bival] (Dbiv) at (3.5, -4.1) {$D_{\mathrm{biv}}$ (bivalent)};

\draw[dashed, gray, rounded corners=5pt] (-4.8,-1.5) rectangle (2.1,-4.8);
\node[font=\small\itshape, gray!60!black] at (-1.35,-1.25)
      {$\mathcal{E}$ (reachable without $e$)};

\draw[arrow, dashed] (C.south west)
    -- node[left, font=\scriptsize, pos=0.55]{avoid $e$} (C0.north);
\draw[arrow, dashed] (C.south east)
    -- node[right, font=\scriptsize, pos=0.55]{avoid $e$} (C1.north);
\draw[arrow, dashed] (C0.south) -- (B.north west);
\draw[arrow, dashed] (B.north east) -- (C1.south);
\draw[arrow] (B.east) -- node[above, font=\scriptsize, pos=0.5]{apply $e=(p,m)$} (Dbiv.west);

\node[draw=blue!55, fill=blue!4, rounded corners=3pt,
      font=\scriptsize, align=center, text width=4.4cm, inner sep=5pt]
      at (-2.6, -6.0)
      {\textbf{Case A:} $e_0$ affects a different process.\
       The events commute, so opposite valencies\
       would reach the same configuration.};

\node[draw=teal!65, fill=teal!4, rounded corners=3pt,
      font=\scriptsize, align=center, text width=4.4cm, inner sep=5pt]
      at (3.3, -6.0)
      {\textbf{Case B:} $e_0$ affects the same process $p$.\
       Deliver $m,m_0$ in both orders, crash $p$,\
       and use indistinguishability to force one decision.};
\end{tikzpicture}
\caption{\textbf{Step 2: Bivalency is preserved after any event.} Let $C$ be bivalent and $e=(p,m)$ any pending event. $\mathcal{E}$ is the set of configurations reachable without applying $e$, and $\mathcal{D}=e(\mathcal{E})$ is obtained by applying $e$ to each of them. Under the contradiction assumption that every configuration in $\mathcal{D}$ is univalent, a path in $\mathcal{E}$ yields a boundary pair with opposite $e$-successor valencies. The two cases above give the contradiction.}
\end{subfigure}
\caption{Proof structure of the FLP theorem. \textcolor{orange!80!black}{Yellow/orange}: bivalent configurations. \textcolor{blue!60!black}{Blue}: 0-valent. \textcolor{red!60!black}{Red}: 1-valent.}
\label{fig:flp-proof}
\end{figure}

\subsection{What FLP Does Not Say}

FLP does not say consensus is impossible in all settings. It applies specifically to \emph{deterministic} algorithms in \emph{fully asynchronous} systems. Practical systems circumvent the theorem in 3 ways:\\[-0.8cm]
\begin{enumerate}[label=(\roman*)]
    \item \textbf{Randomisation.} Algorithms such as Ben-Or's use coin flips to break symmetry, achieving consensus with probability~1 in an asynchronous model. Because the algorithm is no longer deterministic, the FLP argument does not apply.
    \item \textbf{Partial synchrony.} If messages are eventually delivered within some (unknown) time bound, correct algorithms exist. This is the assumption under which Paxos operates.
    \item \textbf{Failure detectors.} Abstract oracles that provide (possibly imperfect) information about crashed processes. Chandra and Toueg~\cite{chandra1996} studied which failure-detector classes suffice for consensus, while Chandra, Hadzilacos, and Toueg~\cite{cht1996} characterised the weakest failure detector for solving consensus in asynchronous systems with crash failures.
\end{enumerate}


% ══════════════════════════════════════════════════════════
% SECTION 5  --  PAXOS: THE PRACTICAL BYPASS
% ══════════════════════════════════════════════════════════
